\section{AI 使用声明}\label{sec:ai-decl}

本节如实说明生成式 AI 在本项目中的使用范围、边界与责任归属。
依据是赛题补充说明（一）允许使用 AI，但应体现参赛者的设计判断，
以及评分标准中论文/技术报告维度对AI 使用声明严谨程度的考察要求。
权威版本见仓库 \texttt{docs/declaration/ai-use-declaration.md}，本节为其正文对照。

\subsection{使用范围}

生成式 AI 在本项目中承担\textbf{辅助}性工作：资料检索线索整理、建模方案讨论、
代码草拟与调试、测试用例设计、文字整理、可视化方案构思，以及\textbf{装饰性}视觉素材的生成。

\subsection{边界}

\begin{itemize}
  \item \textbf{不生成实验数据}：本文中的实验数据、模型指标、统计图与网络结构图
        均由实际程序运行产生，凭证为 \texttt{results/} 下可重生成的产物（见 \S\ref{sec:repro}）。
        AI 生成的装饰性图像\textbf{不作为}实验数据或科学证据。
  \item \textbf{不填未定数值}：按仓库硬约束，上游文档未定义的数值、公式与语义，
        实现侧不得自行选定并当作既定事实。本文引用的每个数字均可回溯到
        \texttt{results/} 或 \texttt{configs/}；未跑完的实验留空并注明，不以估算替代。
  \item \textbf{不替代核心设计}：问题定义、数学模型（RGCD 发育映射、DanioNet 动力学、
        遗传边界与当代学习）、实验设计与结论判断均由团队成员完成；
        AI 仅作为实现与整理的加速工具。
\end{itemize}

\subsection{责任归属}

团队成员对问题定义、数学模型、代码、数据、实验、统计分析、文献引用、
技术报告与最终路演陈述进行\textbf{人工审核}，并\textbf{承担最终责任}。
